\documentclass[11pt]{book}
\usepackage{fontspec}
\setmainfont{Latin Modern Roman}
\setsansfont{Latin Modern Sans}
\usepackage{amsmath,amssymb,mathtools}
\usepackage{unicode-math}
\setmathfont{Latin Modern Math}
\usepackage{xurl}
\usepackage[hidelinks]{hyperref}
\newcommand{\sref}[2]{\hyperlink{source:#1:#2}{[S#2]}}
\newcommand{\cataloguescope}[1]{\paragraph{Mathematical question.}#1}
\begin{document}
\section{Research Problems in Function Theory — Problem 3.30}\label{statement}

\subsection{Definitions and mathematical statement}

\paragraph{Shared notation.}$\mathbb N=\{1,2,\ldots\}$, and $\mathbb Z,\mathbb Q,\mathbb R,\mathbb C$ denote the integers, rationals, reals and complex numbers. Any other conventions are specified locally.


Let $\Gamma\subset\mathbb C\cong\mathbb R^2$ be a smooth Jordan curve, meaning an embedded smooth closed curve bounding a bounded interior region. Let $n(z)$ be its inward unit normal and $ds_z$ arc-length measure. For distinct $z,z'\in\Gamma$, define the double-layer kernel
\[K(z,z')=\frac{n(z)\cdot(z'-z)}{|z'-z|^2},\]
which equals $\cos\phi/|z-z'|$ with the source's inward-normal convention. Use its continuous diagonal extension, which exists for a smooth curve. Define
\[(T_\Gamma f)(z')=\int_\Gamma f(z)K(z,z')\,ds_z\qquad(f\in C(\Gamma)).\]
Its rank is the dimension of its range in $C(\Gamma)$; a circle gives a constant kernel and rank one.
\paragraph{Statement recorded in the supplied source.}
Is there any smooth Jordan curve $\Gamma$ that is not a circle for which $\operatorname{rank}T_\Gamma<\infty$? Equivalently, are circles the only smooth Jordan curves giving finite rank? \sref{statement}{2}
\subsection{Short English statement}
Are circles the only smooth simple closed curves whose double-layer boundary operator has finite rank?
\subsection{Sources}
\begin{description}
\item[\hypertarget{source:statement:1}{S1}] ULAM.AI and the UnsolvedMath Contributors, \emph{UnsolvedMath: A Curated Collection of Open Mathematics Problems}, catalogue entry. CC BY 4.0. \url{https://huggingface.co/datasets/ulamai/UnsolvedMath}.
\item[\hypertarget{source:statement:2}{S2}] Walter K. Hayman and Edward F. Lingham, \emph{Research Problems in Function Theory}, arXiv:1809.07200, 2018. Problem 3.30, its complete Update 3.30, and the relevant chapter notation. \url{https://arxiv.org/pdf/1809.07200}.
\end{description}
\label{end:statement}
\end{document}
